\clearpage
\chapter*{2026 Enero}
\addcontentsline{toc}{chapter}{2026 Enero}

\begin{enunciado}
    \hypertarget{2026 A 1}{}
    \textbf{1. (3 puntos)} Demostrar que un espacio topológico es compacto si y solo si se verifica que en él toda base de filtro \(\{B_\lambda\}_{\lambda \in L}\) tiene un punto de aglomeración.
\end{enunciado}

\begin{enunciado}
    \hypertarget{2026 A 2}{}
    \textbf{2. (3 puntos)} Consideramos el espacio topológico \((\mathbb{R}, T_f)\), siendo \(\mathbb{R}\) el conjunto de los números reales y \(T_f\) la topología final para la aplicación \(f : (\mathbb{R}, T_u) \to \mathbb{R}\), definida, para cada \(x \in \mathbb{R}\), mediante \(f(x) = |x|\) (valor absoluto de \(x\)). Denotamos \(T_u\) la topología usual de la recta real. Encontrar la adherencia y el interior del conjunto \(B = [-1, 1]\) en \((\mathbb{R}, T_f)\).
\end{enunciado}

\begin{enunciado}
    \hypertarget{2026 A 3}{}
    \textbf{3. (4 puntos)} En el conjunto \(\mathbb{R}\) de los números reales, definimos la topología \(T\) mediante \(T = \{\emptyset\} \cup \{A \subset \mathbb{R} \mid \mathbb{R} - \mathbb{Q} \subset A\}\) donde \(\mathbb{Q}\) son los números racionales. Estudiar
    \begin{enumerate}[label=\alph*)]
        \item Si el espacio \((\mathbb{R}, T)\) es \(T_2\).
        \item Si el espacio \((\mathbb{R}, T)\) es regular.
        \item Si el espacio \((\mathbb{R}, T)\) verifica el segundo axioma de numerabilidad.
    \end{enumerate}
\end{enunciado}

